\section{Inherited compiler contracts and the consistency endpoint}
\label{sec:inherited-compiler-contracts}

The new projection and factorial argument are composed with the following
specific arithmetic interfaces.  All variables in this section range over
$\mathbb N=\{0,1,\ldots\}$ unless indicated otherwise.  For the unrestricted
integer equation, the sign-recovery certificate must first establish this
natural domain.  The statements below concern every satisfying assignment;
their use does not presume the intended values of intermediate witnesses.

\subsection{Logical calibration and what is imported}

The numerical source is the fixed twelve-branch calculus $C_{222}$, with
truth code $0$ and falsity code $1$.  Its calibration is
\[
 C_{222}(1)
 \ \longleftrightarrow\
 \operatorname{BProv}(\mathrm{ZF},\bot)
 \ \longleftrightarrow\
 \operatorname{BProv}(\mathrm{ZFC},\bot).
\]
The exact named results are
\nolinkurl{KEndpoint.c222_iff_zf_refutation} and
\nolinkurl{ConstructibleChoice.zf_refutation_iff_zfc_refutation}.
Their frozen files, relative to the direct-equation component, are
\nolinkurl{provenance/upstream/calibration/KEndpoint.lean} and
\nolinkurl{provenance/upstream/calibration/ConstructibleChoice.lean}.
The source numerals and rules are in the adjacent
\nolinkurl{Code254.lean} and \nolinkurl{KCode222.lean}.
The repository paths are
\nolinkurl{Computability/BusyBeaver/Lean/BusyBeaver/ZFC/}, at calibration
commit \nolinkurl{08200eb10fde2db408c69dffd840d6920a569b1f}; exact file hashes
are retained in the provenance ledgers.

These are substantial imported mathematical results.  The first combines
the soundness of the numerical calculus with its model-existence
completeness argument; the second removes Choice by the constructible-model
argument and first-order soundness/completeness.  No Turing-machine endpoint
is used.  The appeal is to these theorems and their proofs formalizable in
ZFC, not to an inference that every external Lean theorem is automatically
a theorem of ZFC.  The complete transitive Lean development was not rebuilt
for this contribution, and no new object-language ZFC proof term is claimed.

The remaining passage from finite rule instances to a proof is the
single-conclusion Horn-flow criterion.  Give $(p,q)\to v$ the vector
$e_v-e_p-e_q$, with $e_0=0$.  A proof tree has total flow $e_1$.
Conversely, if a finite nonnegative sum of rule vectors equals $e_1$,
sum coordinates outside the least rule-closed set containing $0$.
Each rule contributes at most zero outside that set; a missing target
would contribute one, a contradiction.  Cycles do not invalidate this
argument.

For the accumulator-to-flow step, all restored labels lie in
$[0,L]$, where $L=8m^2+8m+2$, and $M=(L+1)^t$, $t\ge1$.
Telescoping the accumulator equations with endpoints $S_0=1+M$, $S_t=1$
gives $P(M)=Q(M)$, where
\[
 P(X)=\prod_{i<t}(1+v_iX),\qquad
 Q(X)=(1+X)\prod_{i<t}(1+p_iX)(1+q_iX).
\]
Since $P(M)<M^{t+1}$, nonnegativity first forces $\deg Q\le t$.
Both products then have at most $t$ nonconstant factors, so every
coefficient is strictly below $M$.  The least nonzero coefficient of
$P-Q$ could not vanish on evaluation at $M$; hence $P=Q$.
Unique factorization in $\mathbb Q[X]$ gives precisely the required flow.
This is the bounded-factor argument in
\nolinkurl{provenance/upstream/calibration/direct_flow_large_carrier_README.md}.
Conversely, expand a proof tree and let $S_i$ be the product of the weights
of its open obligations.  A finite CRT beta code gives
$x=S_i+(1+(i+1)y)u_i$, choosing $y$ larger than every state and divisible
by $1,\ldots,t+1$.  In soundness, the supplied bounds
$S_i\le y<1+(i+1)y$ make these actual remainders and make adjacent
rows agree.  Thus this encoding is an equivalence of existence predicates.

\subsection{The bounded-row CRT contract}

Let $E,F$ be polynomials with nonnegative integer coefficients in the row
index, fixed outer parameters, and $r$ local coordinates.  In the present
source $r=9$ and the coordinate bounds are
\[
 \mathbf u=(m,m,m,m,m,y,y,x,x).
\]
Suppose $t\ge1$ and
\[
 H\ge E(t,\mathbf u)+F(t,\mathbf u)+t+x+y+m+1.
\]
With the outer parameters suppressed, put
$T=H!$, $d_i=1+(i+1)T$, and
$\mathcal B=\prod_{i<t}d_i$.  Then the bounded family
\[
 \forall i<t\ \exists w\le\mathbf u:\ E(i,w)=F(i,w)
\]
is equivalent to the existence of natural $C,A_1,\ldots,A_r$ such that
\begin{align*}
 \mathcal B&\mid\binom{A_j}{u_j+1} &&(1\le j\le r),\\
 \mathcal B&\mid 1+(C+1)T, &
 E(C,A)&\equiv F(C,A)\pmod{\mathcal B}.
\end{align*}
The representatives $A_j$ may all be required positive, as they are in the
literal compiler.

For soundness, choose a prime $\ell\mid d_i$.  Since $d_i$ is coprime to
$H!$, necessarily $\ell>H$.  Multiplying the range congruence by
$(u_j+1)!$ gives a zero falling factorial modulo $\ell$, so
$A_j\equiv w_j\pmod\ell$ for some $0\le w_j\le u_j$.
The index congruence gives $C\equiv i\pmod\ell$ by cancellation of $T$.
The two evaluations are congruent and both lie in $[0,H]$, hence are equal.
For completeness, the $d_i$ are pairwise coprime: a common prime would
divide $j-i$ and so divide $H!$, which is impossible.  CRT encodes the
chosen row witnesses.  Each $(u_j+1)!$ divides $H!$ and is coprime to
$\mathcal B$, so it cancels from the falling-factorial congruence.
Positive representatives follow by adding multiples of $\mathcal B$.
The literal evaluation congruence uses
$E+\mathcal B q=F+\mathcal B v$ with natural $q,v$; arbitrary signed
quotients are represented by choosing the common offset $v$ sufficiently
large, so this orientation adds no restriction.

The supplied majorant certificate checks 392 gate inequalities and bounds
$E+F+t+x+y+m+1$ by $R_0^{6,893,174}$, where
$R_0=2+t+x+y+m+M_0\ge2$.  The literal choice
$H=R_0^{2^{23}}$ is therefore sufficient.  Each addition uses the exponent
bound $\max(e_1,e_2)+1$, each multiplication uses $e_1+e_2$; induction on
the actual nonnegative pair circuit proves the bound.  The factorial
bootstrap supplies the genuine value of $T$ before this CRT theorem is
invoked.

The fixed-length modulus product is itself represented arithmetically.
For $Q=(1+(t+1)T)^t+2$ and $D=TQ-1>0$,
\[
 \mathcal B=
 \left(T^t\,t!\binom{Q+t}{t}\right)\bmod D.
\]
Indeed the numerator is $\prod_{i=1}^t(TQ+iT)$, congruent to
$\prod_{i=1}^t(1+iT)$ modulo $D$.  The latter product is strictly smaller
than $D$ by the defining bound on $Q$.  An ordinary quotient/remainder
equation with remainder slack supplies this graph; no unbounded product
operator remains in the polynomial.

\subsection{One radix for all binomial requests}

For the twelve actual requests $\binom{n_j}{k_j}$, all $n_j,k_j$ are
positive.  Let $N$ be the sum of the distinct upper-argument circuit atoms
and $q=2^N+1$.  Thus $q>2^{n_j}$.  With
$L_j=q^{k_j}$ and $U_j=(q+1)^{n_j}$, the literal natural equations are
\[
 U_j=L_j(qh_j+C_j)+r_j,\qquad
 L_j=r_j+s_j+1,\qquad q=C_j+d_j+1.
\]
They recover exactly digit $k_j$ of the carry-free base-$q$ expansion of
$(q+1)^{n_j}$; therefore $C_j=\binom{n_j}{k_j}$, including $k_j>n_j$.
Conversely ordinary division supplies all four natural auxiliary values.
This theorem is simultaneous for arbitrary assigned argument values.
In completeness, the shared radix is selected after all semantic upper
arguments, including the chosen CRT representatives, have been fixed.

\subsection{Shared Pell certificates for all positive powers}

Write
$(a+\sqrt{a^2-1})^j=X_j(a)+Y_j(a)\sqrt{a^2-1}$ for $a\ge2$.
The imported number-theoretic ingredients are classification of natural
Pell solutions and the finite Matiyasevich index criterion.  The exact
repository interface is
\nolinkurl{DiophantineConsistency.PellRepresentation.criterion_iff_pell},
in \nolinkurl{provenance/upstream/calibration/PellRepresentation.lean},
which invokes \nolinkurl{Pell.matiyasevic} from Mathlib's
\nolinkurl{Mathlib/NumberTheory/PellMatiyasevic.lean}.

For the 23 distinct power requests $(n,k,m)$, the common frame is
\[
 W=2+\sum\{\hbox{distinct base, exponent, and output atoms}\},\quad
 Z=z+1,\quad A^2-W^3(W+2)Z^2=1.
\]
Each of the seven distinct bases has a natural $T_n$ satisfying
$2An=T_n+n^2+1$.  This equation excludes $n=0$ before any factorial or
power value is recovered.  Every exponent is syntactically positive.
Pell classification applied at base $W+1$ gives
$A=X_j(W+1)$ and $WZ=Y_j(W+1)$.
Modulo $W$, the recurrences give $Y_j(W+1)\equiv j$; thus
$j\ge W$.  Since $1\le n,k\le W<A$ and $m<W$, consequently
\[
 A\ge(W+1)^W>W^W,\qquad
 T_n\ge2A-2>W^W\ge n^k,\qquad m<T_n.
\]
Conversely $j=W$ supplies a frame for any finite family of intended
values.  The weaker inequality $A>W^2$ alone is insufficient for power
recovery; the divisibility of the frame's Pell index is essential.

For each of the eighteen exponent atoms $k_i$, use natural
$X_i,Y_i,S_i,R_i,e_i,s_i,r_i$.  The literal shared index equations are
\begin{align*}
 Y_i&=k_i+e_i,&X_i^2-(A^2-1)Y_i^2&=1,\\
 S_i^2-(B^2-1)R_i^2&=1,&S_i&=X_i+Us_i,\\
 R_i&=k_i+4Y_ir_i,&U^2-(A^2-1)V^2&=1,\\
 V&=\left(\prod_iY_i^2\right)(v_0+1),&
 B&=1+4V(b_0+1)=A+Ub_1.
\end{align*}
They imply every hypothesis of the nonzero branch of the imported index
criterion: $A,B>1$, $Y_i\ge k_i>0$, $V>0$, $Y_i^2\mid V$, and the four
required congruences.  Hence $(X_i,Y_i)=(X_{k_i}(A),Y_{k_i}(A))$.
For simultaneous completeness, set $L=4\prod_iY_i^2$ and choose a positive
power of the unit $A+s$ equal to one in the finite ring
$(\mathbb Z/L\mathbb Z)[s]/(s^2-(A^2-1))$.
Its Pell coordinates give $U\equiv1\pmod L$, $V\equiv0\pmod L$, $V>0$,
and $\gcd(U,4V)=1$.  CRT chooses $B>A$ with the displayed congruences.
Set $S_i=X_{k_i}(B)$, $R_i=Y_{k_i}(B)$; recurrence congruences and
monotonicity give all required natural quotients.

Finally each request has one natural quotient $h$ and
\[
 X_i+nY_i=AY_i+m+T_nh.
\]
For $H_j=X_j(A)-(A-n)Y_j(A)$, the recurrence is
$H_{j+1}=nH_j+T_nY_j(A)$, $H_0=1$.  Thus
\[
 H_k-n^k=T_n\sum_{j<k}n^{k-1-j}Y_j(A).
\]
Both $m$ and $n^k$ are in $[0,T_n)$, proving soundness; the displayed
nonnegative sum supplies $h$ for completeness.  The finite-ring and
recurrence arguments justify shared witnesses; merely identifying
independently existential variables would not suffice.

\paragraph{Composition inside ZFC.}
These proofs use finite induction, integer/rational arithmetic, prime
factorization, CRT, finite groups, the stated Pell theorems, and the
identified logical calibration.  They can be carried out over ZFC's own
$\omega$ and coded finite sequences, rather than only for externally
standard numerals.  After the factorial fixed-point proof and exact
polynomial projections are composed with these contracts, ZFC proves
existence of a polynomial zero equivalent to a coded ZFC refutation;
negating gives the asserted consistency biconditional.  Exact file checks
verify the supplied finite circuits and coefficients.  They do not prove
consistency or replace any imported semantic theorem.
